\input{../tex-inputs/latex-leseansicht-vorspann}

\section[Arthur und Olga Schnitzler an Gustav Schwarzkopf, 19. 9. 1910]{L04370 Arthur und Olga Schnitzler an Gustav Schwarzkopf, 19. 9. 1910}
\nopagebreak\mylabel{L04370v}
\rehead{ }\normalsize\beginnumbering\briefempfaengerindex{Schwarzkopf, Gustav@\textsc{Schwarzkopf, Gustav}!zzzSchnitzler, Olga@\emph{von Olga Schnitzler}!1910-09-191@{19. 9. 1910}|(be}\briefempfaengerindex{Schwarzkopf, Gustav@\textsc{Schwarzkopf, Gustav}!zzzSchnitzler, Arthur@\emph{von Arthur Schnitzler}!1910-09-191@{19. 9. 1910}|(be}
\toendnotes[C]{\smallbreak\pagebreak[2]}
\correspDesc{Versand  durch Arthur Schnitzler, Olga Schnitzler am 19. 9. 1910 in Frankfurt am Main
\newline{}Übermittlung  am 19. 9. 1910 in Frankfurt am Main
\newline{}Erhalt  durch Gustav Schwarzkopf im Zeitraum [20. 9. 1910 – 24. 9. 1910?] in Wien}\toendnotes[C]{\smallbreak}
\Standort{DLA, A:Schnitzler, HS.1985.1.1897.}
\physDesc{Bildpostkarte, 174 Zeichen
\newline{}Handschrift Arthur Schnitzler: Bleistift, deutsche Kurrent
\newline{}Handschrift Olga Schnitzler: Bleistift, lateinische Kurrent
\newline{}Versand: Stempel: »\nobreak{}\oindex{Frankfurt am Main@\textbf{Frankfurt am Main}|pwk}Frankfurt (Main) \textcolor{gray}{9}, 19. 9. 10, 5–6N\nobreak{}«.  }\toendnotes[C]{\smallbreak}\pstart{}{\pb}\textsc{Herrn Gustav Schwarzkopf}\pend{}\pstart{}Wien I\oindex{XXXX Ortsangabe fehlt|pw}\pend{}\pstart{}\textsc{Tiefer Graben 23}\oindex{Wien@\textbf{Wien}!I., Innere Stadt@\textbf{I., Innere Stadt}!Tiefer Graben 23@\textbf{Tiefer Graben 23}|pw}.\pend{}{\bigskip}
\pstart
           \noindent{}\centering{}{\pb}\textcolor{gray}{\textbf{Frankfurt a. M.}}\oindex{Frankfurt am Main@\textbf{Frankfurt am Main}|pw}\pend
           
\pstart
           \centering{}\textcolor{gray}{\textbf{Goethehaus}}\oindex{Goethe-Haus@\textbf{Goethe-Haus}|pw}\pend
           
\pstart
           \centering{}\textcolor{gray}{\textbf{Gr. Hirschgraben 23}}\oindex{Goethe-Haus@\textbf{Goethe-Haus}|pw}\pend
           
\pstart
           \centering{}\textcolor{gray}{\textbf{(16. Jahrh.)}}\pend
           \vspace{1em}
\pstart
           \raggedleft{}{\pb}19. 9. 1910.\pend
           \vspace{0.5em}
\pstart
           Herzliche Grüße u auf
      baldgs Wiederſehen. Die \label{K_L04370-1v}\edtext{Oper\pwindex{Schnitzler, Arthur 15. 5. 1862 Wien – 21. 10. 1931 ebd.@\textsc{Schnitzler, Arthur} (15. 5. 1862 Wien – 21. 10. 1931 ebd.)!Liebelei. Oper in drei Akten@\strich\emph{Liebelei. Oper in drei Akten}|pwv}\eventindex{Frankfurt am Main@\textbf{Frankfurt am Main}!Uraufführung von Liebelei. Oper, 18.9.1910@Uraufführung von Liebelei. Oper, 18.9.1910|pwv} war ein großer
               Erfolg}{\lemma{\textnormal{\emph{Oper … Erfolg}}}\Cendnote{\textnormal{Vgl. Uraufführung von Liebelei. Oper, 18.9.1910. In: A. S.: \emph{Kulturveranstaltungen}, \url{https://schnitzler-kultur.acdh.oeaw.ac.at/pmb46947.html}.}}}\label{K_L04370-1}; ob auch weit,
      wird die Zukunft
      lehren.\pend
           
\pstart
           Ihr{\\[\baselineskip]}\spacefill\mbox{A.}\pend
           \leftskip=0em{}\selectlanguage{ngerman}\endnumbering\briefempfaengerindex{Schwarzkopf, Gustav@\textsc{Schwarzkopf, Gustav}!zzzSchnitzler, Olga@\emph{von Olga Schnitzler}!1910-09-191@{19. 9. 1910}|)be}\briefempfaengerindex{Schwarzkopf, Gustav@\textsc{Schwarzkopf, Gustav}!zzzSchnitzler, Arthur@\emph{von Arthur Schnitzler}!1910-09-191@{19. 9. 1910}|)be}\mylabel{L04370h}
\begin{anhang}
\end{anhang}\newcommand{\dateiname}{L04370}\newcommand{\titel}{Arthur und Olga Schnitzler an Gustav Schwarzkopf, 19. 9. 1910}\newcommand{\editorInnen}{Herausgegeben von Jahnke, SelmaMüller, Martin Anton}\input{../tex-inputs/latex-leseansicht-abspann}
